% =============================================================================
% Chapter 11: Changelog
% =============================================================================

\chapter{Changelog}
\label{ch:changelog}

This changelog follows the Keep a Changelog format~\citep{keepachangelog}.
Versions are listed in reverse chronological order. This chapter tracks the
\emph{document} baseline; library releases carry their own
\code{CHANGELOG.md} in the repository (REQ-94). From the next document
version onward, every entry lists the requirement(s) and design decision(s)
it adds or modifies.

\section*{Document 0.2.4, 2026-07-30, When a Declared Uncertainty Applies}

A PATCH revision by chapter~9's rule: a clarification, adding no
requirement and withdrawing none.

\paragraph{Changed (Section~\ref{sec:itceq}, the \code{[uncertainties]}
application moment).} The section rules now state WHEN a declared
uncertainty is applied: before the first equation for a name the file does
not produce, and immediately after the writing equation for a name it
does. The rule keys on what the file produces and never on what the
VarFrame carried on the way in.

The document had said only that the values are absolute or relative and
per-variable, so the ordering was unspecified and the implementation had
adopted the other rule: it sorted by whether the incoming frame already
carried the name. A target the frame carried was therefore assigned before
the loop and then overwritten by its own equation's propagation. With
\code{[uncertainties] y = 5.0} alone, the declaration was lost and the
frame shipped no uncertainty at all; with \code{x = 1.0} beside it, the
same file shipped \code{u(y) = 2.0}, the value \code{u(x) = 1.0}
propagates to through \code{y = "2*x"}, where the file declares
\code{5.0}. A relative declaration was resolved against the stale input
column rather than against what the equation wrote.

The wrong-value branch is why this is stated in the specification and not
only fixed in code. A lost declaration is a visible gap a reader can
notice. A different, finite, plausible uncertainty, chosen by whether the
input frame happened to carry a column, is a number no reader can check,
which is this library's worst failure mode
(\code{R4-ITA-003}, \code{REV-004}, \code{ITC-20260730-0105}).

\section*{Document 0.2.3, 2026-07-30, The Sixth Refusal and the Build}

A PATCH revision by chapter~9's rule: clarifications only, no requirement
added and none withdrawn.

\paragraph{Changed (REQ-01, the sixth CHK-1 refusal).} \code{itc.load}
refuses a CSV row WIDER than its header with \code{DataError}, naming the
row, its field count, the header, the cells that would have been
discarded, and the suggested fix. Rows are read by header position, so
such a cell would be dropped with no signal while Provenance went on
hashing the complete file bytes, and the source hash would certify content
the frame did not represent (\code{R3-ITA-009}). REQ-01 now also states
the ASYMMETRY: a row NARROWER than its header is not refused, its absent
trailing cells being NaN-filled exactly as an empty cell is, and the
empty-cell rule itself is stated because the asymmetry leans on it and the
document gave it nowhere. The asymmetry is deliberate, a too-wide row
being a malformed file and a too-narrow one ordinary missing data, and a
reader who meets one half infers the other wrongly.

This refusal shipped in v0.2.0, marked BREAKING in \code{CHANGELOG.md},
with its own error message citing REQ-01 while REQ-01 did not describe it.
The 0.2.2 revision below closed five gaps of exactly this shape and missed
this one; the independent review REV-004 found it
(\code{ITC-20260729-2255}). So the CHK-1 checkpoint added SIX public
refusals in total, five closed there and one here.

\code{CHANGELOG.md} counts FIVE, and both numbers are right about
different things. It counts the rules a caller can hit and folds
\code{validate}'s third refusal into the constant-shadowing entry it
shares a cause with, while this document states that one as its own
amendment to REQ-45. Recorded in both places because a reader comparing
them otherwise concludes one is missing a refusal.

\paragraph{Changed (REQ-95, the specification build).} REQ-95's CI
enumeration gains a job that builds this document and fails on a non-zero
exit or a logged error. Nothing compiled the specification before, so a
document that could not be built was indistinguishable from one that
could. Two defects had already reached \code{main} that way: chapter~7
carried a blank line after every content line, making every source line
its own LaTeX paragraph and splitting a \code{caption} across a paragraph
break, and one unescaped underscore inside a code macro argument put chapter~8 into
math mode that leaked into chapter~9. Seven \code{pdflatex} errors between
them, on a document whose \code{main.log} was dated 2026-07-23 while
document versions 0.2.1 and 0.2.2 both shipped after it
(\code{ITC-20260729-2300}, \code{ITC-20260729-2350},
\code{ITC-20260730-0010}).

The chapter~7 reflow is recorded as a NON-CHANGE to content, which is the
part worth stating: 692 source lines became 346, and the 307 non-blank
lines are byte-identical and in the same order. REQ-96's pre-commit mirror
deliberately does NOT gain the build, since a TeX toolchain is not a
pre-commit dependency; that deviation from REQ-96's mirror obligation is
stated in the requirement rather than left to be noticed.

\section*{Document 0.2.2, 2026-07-29, CHK-1 Refusal Amendment}

Issued to close \code{ITC-20260729-1450}, which blocked the v0.2.0 tag.
The CHK-1 release checkpoint added five public refusals to surfaces
whose requirement text did not describe them, which is the defect class
the 0.2.1 revision above named in its own words and did not finish
closing. The code shipped correct; the text was wrong. No requirement
scope is added or withdrawn here.

A SIXTH refusal from the same checkpoint was missed by this revision and
is closed by 0.2.3 below. The count here is left at five rather than
corrected upward, because five is what this revision did and a revision
entry is a record of one.

\paragraph{Changed (the five CHK-1 refusals).} REQ-01 states that one
name may not be both a dimension and a variable, so a source column
named \code{datapoint} is refused at load and at VarFrame construction
(\code{R3-ITA-008}). REQ-24 and REQ-107 both state that
\code{itc.concat} raises \code{AxesError} when the inputs express a
vector group in different axis systems, or declare one group over
different components, with an undeclared input counting as the
canonical body axis (\code{CHK1-003}). REQ-39 states that a declared
uncertainty magnitude must be finite, with the limit of that check
(declared magnitudes only, OQ-40) stated rather than implied
(\code{R3-ITA-007}, \code{ITACA-031}). REQ-44 states that a built-in
expression constant may not shadow a measured channel, and REQ-45
raises \code{validate} from two refusals to three with the same
collision on the \code{.itceq} path (DD-42, \code{CHK1-001}, OQ-41).

\paragraph{Changed (REQ-92, a false enumeration).} REQ-92 claimed that
\code{0.2.0} ships three breaking changes, and it ships more; the CHK-1
checkpoint alone added several after the carve-out was written. The
remedy is not a corrected count inside the requirement, because a bare
count is the same stale artifact with its evidence removed. The
enumeration and the count are both removed, and REQ-92 now points at
\code{CHANGELOG.md}, which the requirement's own next paragraph already
mandates as the place every breaking change is marked with its
migration. A second inventory in a normative document goes stale
independently of the one it duplicates, and this one did.

\paragraph{Changed (REQ-98 and REQ-26, one count in place of three).}
REQ-98's provisional family was carried in four places at three
different values, four and three inside the requirement itself and two
in the workspace charter, while the normative
Table~\ref{tab:unc-ops} still stated that \code{smooth} and \code{diff}
propagate through their weights, which the implementation has never
done. The list is now given once IN THIS DOCUMENT, at FIVE (the
workspace charter and the reviewer charters outside it still restate a
subset, which this document cannot close and does not claim to):
\code{smooth}, \code{diff}, \code{fitmodel}, \code{fitvalue}, and
\code{fill} with \code{method="polyfit"}. The table gains the
\code{fill} polyfit row, its \code{smooth}/\code{diff} row is corrected
to raise, and its \code{interpolate}/\code{fill} row is scoped to the
two fill methods that do propagate. REQ-26 carries the same split.

The ground recorded for \code{fill(method="polyfit")} is narrower than
the one first written, and the correction is itself worth recording. It
is not that the fit has data-dependent weights, which was measured false
(superposition departs by $1.3 \times 10^{-13}$) and would have drawn
\code{interpolate} in as well; nor is it that no polynomial weight rule
exists, since one ships and \code{db.interpolate(method="polyfit")}
propagates through it. It is that the FILL path emits no weights for the
propagation to consume. Whether the shipped matrix is admissible over a
support chosen per gap from the non-NaN cells is OQ-42, opened in this
revision, and answering it may return the count to four.

\paragraph{Changed (GUM citation).} REQ-39 and the
Chapter~\ref{ch:standards} GUM mapping cite clause~5.1.2 for the
finiteness rule, by its two operative sentences: the positive square
root excludes NaN, and the definition of $u_c(y)$ as an estimated
standard deviation characterizing a dispersion of attributable values
excludes an infinity. The requirement records that neither exclusion is
inferred from the first-order expansion of clause~5.1, which requires
finiteness of the input $u(x_i)$ nowhere, and that the only domain
restriction in 5.1 is 5.1.6, $|y| \neq 0$, which governs relative
uncertainty.

\section*{Document 0.2.1, 2026-07-28, Independent-Review Amendment}

\paragraph{Changed.} REQ-103 is
restated as a semantic guarantee: two VarFrames in the same semantic
state have the same hash, with same semantic state defined
representationally at the floating-point boundary (DD-40). Its scope
gains dimension and variable metadata and names the axis registry. The
practical consequence is that unit is inside the guarantee even though
Section~\ref{sec:dimension-object} calls it metadata, because
\code{rotate} interprets it. REQ-107 moves from draft to stable and
records that \code{db.rotate} updates a group's axis. REQ-36's guard is
stated as per-subtree. REQ-45 records the \code{[constants]} collision
refusal (DD-39). Requirement scope is otherwise unchanged.

\paragraph{Changed (requirements amended by the same review).} REQ-40
gains the positive-semidefinite condition on the assembled matrix: the
pairwise bound is necessary and not sufficient. REQ-38 gains the
correlation write-back from the transformed covariance and its three
refusal conditions. REQ-44 states that calls take positional arguments
only and that a keyword raises. REQ-71 and P-02 declare
\code{db.to\_pandas()} and \code{db.to\_numpy()} deliberately lossy
interoperability bridges, not covered by the every-export rule, with
the draft-mode warning still applying to them as a safety property.
REQ-36's guard is restated as per-subtree. REQ-45 records the
\code{[constants]} collision refusal (DD-39). REQ-103 is restated as a
semantic guarantee (DD-40) and REQ-107 moves from draft to stable.

\paragraph{Corrected (false or overstated claims).} Chapter~8 stated
that v0.1.0 supported PROV-N and PROV-JSON export; no such symbol
exists and it is an M2 deliverable. Chapter~8 also stated that linting
enforced the NumPy docstring sections; ruff enforces presence and
style, and section completeness on the declared public surface is
enforced by a test, with the Raises and Examples halves not yet
enforced by anything. Fifteen occurrences of ``RPN'' across seven
chapters became ``AST'', the canonical term for the expression engine,
and the now-unused acronym entry was removed.

\section*{Document 0.2.0, 2026-07-23, M1 Re-Baseline and Library Infrastructure}

Issued at the M1 kickoff (ultraplan Batch A approval), together with
the approved execution plan committed as
\code{docs/M1\_EXECUTION\_PLAN.md}.

\paragraph{Added.} Functional requirements REQ-104 (options registry
with exact keys), REQ-105 (typed no-default sentinel), and REQ-106
(accessor registration), all draft, adopted from the 2026-07-23
library-review triage (decisions D1, D2, D6) in the new
Section~\ref{sec:library-infrastructure}. Design-decision boxes DD-23
(co-development with pyflightstream, recorded in the decision log on
2026-07-23 and entered here at this baseline per the editing rule),
DD-24 (options-registry adoption), and DD-25 (\code{uncertainties} as
a dev-only test oracle, never a runtime dependency) in
Chapter~\ref{ch:architecture}, each with its long-form entry in the
repository decision log.

\paragraph{Changed.} Chapter~\ref{ch:roadmap}: the M1 deliverables
table splits into v0.2.0 (computation: operations, axes, pipeline,
processor infrastructure, REQ-105, REQ-106) and v0.2.1 (presentation
and builtins: REQ-104, plot core, WT builtins, statistics and
compare), recording the 2026-07-23 re-baseline decision; the stretch
scope may still ship inside v0.2.0 if green in the same window.
REQ-101 was validated as written at the Batch A checkpoint; the
reqbox stays draft until the implementation and its tests land, then
moves to stable through this changelog and the revision history.
At the M1 Phase B1 checkpoint, REQ-31 gained a normative note that
\code{fitmodel} raises when uncertainty is present (its
coefficient-space rule is not yet frozen, OQ-24), and the
Table~\ref{tab:unc-ops} normative UncFrame semantics gained a
\code{fitmodel}/\code{fitvalue} provisional row and a scoped
restatement of the linearity claim (weight-based rows exact, Jacobian
rows linearized, provisional rows raising). OQ-24 (coefficient-space
uncertainty) and OQ-25 (origin-tag reduction across collapsing and
expanding operations) were opened during the same phase. Three
author-seat calls were resolved at the B1 checkpoint: \code{fitvalue}
defers with \code{fitmodel} and raises on uncertainty (REQ-32,
OQ-24); the exception hierarchy replaces the \code{diff}-only
\code{DiffWindowError} with a shared \code{FitDegreeError} for the
too-few-points-for-degree invariant across \code{diff}, \code{smooth},
\code{interpolate}, and \code{fitmodel} (REQ-30,
Table~\ref{tab:exception-hierarchy}); and \code{fill}'s positional
\code{method} argument is deprecated toward keyword-only for
consistency with the M1 kernel operations (REQ-26).

\paragraph{Added (B2).} Functional requirement REQ-107 (per-group
source frames and multi-frame data), draft. Design decisions DD-26
(\code{scipy} as a dev-only direction-cosine test oracle) and DD-27
(the sanctioned accessor per-instance cache write) in
Chapter~\ref{ch:architecture}. The axis registry and vector-group
declarations join the state hash (REQ-103). The built-in wind and
stability frames follow the AIAA R-004A Etkin direction-cosine
convention, SME-accepted at the B2 checkpoint. The exception
hierarchy gains \code{AxisNotFoundError}, \code{VectorGroupError},
\code{RotationMatrixError}, and \code{AccessorRegistrationError}.

\paragraph{Changed (B2).} REQ-101 (condition-dependent axes) promoted
from draft to stable: condition-dependent frames are implemented and
tested (per-grid-point evaluation, chain-rule angle uncertainty,
proper-rotation validation, angle unit from the Dimension or Variable
metadata). The rotation propagation treats frame angles as mutually
independent, rejecting a declared correlation that involves an angle
rather than dropping it; the modeling question is recorded as OQ-26.
Three B2 author-seat calls were resolved (queue Q-013): \code{Axis}
is exported at the package top level; the axes surface standardizes on
``axis'' for a coordinate system (REQ-107 \code{declare\_vector(axis=)},
distinct from the \code{Frame=} selector of \code{select}); and
\code{translate\_moments} gains \code{force} and \code{moment} group
selectors (REQ-100). OQ-25 (origin-tag reduction across collapsing and
reshaping operations) was resolved and its four rules folded into the
HistoryFrame section (Chapter~\ref{ch:datamodel}, Section 4.3).

\paragraph{Changed (B3a).} The \code{.itc\_pipe} format
(Chapter~\ref{ch:datamodel}, Section~\ref{sec:itc-pipe}) is specified as JSON rather
than TOML, with design decision DD-28 recording why (no
standard-library TOML writer in any Python version, no TOML null type
against a meaningful \code{compute(fill=None)}, and nested replay
arguments) and the reasoning that every required content item is kept:
the creating version, the source index range, each call with its
keyword arguments and attached comment (REQ-19), and a content hash
that is reverified on load. REQ-53 gained the normative extraction
contract: frame construction (\code{load}, \code{pivot}) is input
preparation and is skipped when it leads the range, any other
operation that records no replayable step raises
\code{PipelineCompatibilityError}, and a range yielding no step raises
rather than returning a pipeline that would apply as a silent no-op.
REQ-54 gained the replay contract: the recorded call and comment
reproduce the state hash when replayed onto the frame the range was
lifted from (onto a different frame the processing is reproduced, not
the hash), and an incompatible target raises with the
underlying error's object and suggested fix preserved. The History-entry
table gained the \code{step} field (replay metadata, excluded from the
state hash, REQ-103 scope unchanged), and the \code{.itc} archive
persists it at schema \code{itaca-itc/2} so a reopened archive can still
lift its recipe. The archive section (Section~\ref{sec:itc-archive}) was
amended to match, naming both schema versions, the \code{step} member of
\code{history.json}, and the \code{steps\_hash} digest in
\code{metadata.json}, with the scope boundary against the REQ-103 state
hash stated there and recorded in DD-30. REQ-103 scope is unchanged.
DD-28 also records the replay representation itself:
structured step re-dispatch rather than re-parsing the History display
strings.

\paragraph{Changed (B3b).} \stable{} REQ-83
(Chapter~\ref{ch:nonfunctional}) is amended: the language baseline in
Table~\ref{tab:dependency-versions} rises from \code{>=3.10,<3.14} to
\code{>=3.11,<3.14}. The reason is REQ-48. \code{.itceq} is specified
as a TOML-structured file (Chapter~\ref{ch:datamodel},
Section~\ref{sec:itceq}), and only from 3.11 does the standard library
ship a TOML reader, so the parser now reads it with \code{tomllib}:
no dependency, no vendored reader, and no hand-written parser. This
resolves OQ-28, open since 2026-07-23 and blocking phase B3b, by
taking none of the three options it offered (a vendored reader, a hand
parser, or \code{tomli} as a conditional dependency with a REQ-82
amendment); DD-32 records why, and why DD-28 and the JSON
\code{.itc\_pipe} encoding stand unchanged, \code{tomllib} being
read-only. REQ-83 gains two normative additions: the baseline may not
fall below a standard-library module the implementation depends on,
and the floor is declared once with its restatements (the PyPI
classifiers, the \code{ruff} target version, and the lowest leg of the
CI test matrix) pinned to that declaration. The table caption now
names the amending checkpoint rather than the v0.1.0 baseline.

\paragraph{Changed (B3b, REQ-82).} \stable{} REQ-82
(Chapter~\ref{ch:nonfunctional}) states the NumPy-only scope as every
package except \code{io/} and \code{utils/}, replacing the list of the
three packages that existed at the 0.1.0 baseline. This is what CI
already enforced, since the \code{ruff} ban is repository-wide with
per-file exemptions; the requirement text had not kept up. Stated as an
exception list, a package added later is restricted by default rather
than inheriting nothing at the moment it is least reviewed, which
\code{pproc/} demonstrated in this same phase by arriving covered with
no text change. DD-33 records the reasoning and the matching defect in
the guard test, which enumerated the same package names and has been
changed to discover them. The DD-02 box in
Chapter~\ref{ch:architecture} gained a pointer; the decision itself is
not overturned and its statement still holds.

\paragraph{Changed (B3b, REQ-46).} \stable{} REQ-46's factory is
renamed from \code{itc.pproc(name\_or\_path, config=\{\})} to
\code{itc.processor(name\_or\_path, config=\{\})}, and the top-level
API surface of Chapter~\ref{ch:architecture} moves with it. This
resolves OQ-29. Binding the factory under the package's own name
shadows the attribute the import machinery sets for \code{itaca.pproc},
so the \code{pproc.statistics} and \code{pproc.compare} forms in which
REQ-49 and REQ-50 are specified could never have resolved, and
\code{import itaca.pproc as pp} would have returned a function
silently. Renaming the factory rather than the package leaves all
thirteen module-path citations and all sixteen \code{pproc.<callable>}
citations true as written; DD-34 records the measurement. The worked
example in Chapter~\ref{ch:contributing} was amended in the same
change, because it used \code{pproc} for a third thing, the constructed
processor instance, and DD-04 now names the factory where it shows the
call shape. Requirement scope is otherwise unchanged: REQ-45, REQ-47
and REQ-48 are untouched.

\paragraph{Added (B3b).} Processor infrastructure, REQ-45 to REQ-48
with DD-16 and DD-17, implemented in the new \code{pproc/} package of
Chapter~\ref{ch:architecture}: the \code{Processor} protocol, the
\code{itc.processor()} factory, the idempotence policy, and the
\code{.itceq} parser with cycle detection at parse time and opt-in
topological sorting that reports the resolved order. No requirement
text changed: the \code{ProcessorError} leaves of
Table~\ref{tab:exception-hierarchy} were specified at the 0.1.0
baseline and are implemented here as written, with
\code{ProcessorIdempotenceWarning} carrying both roles REQ-47 gives it,
raised when a reapplication is refused and warned when \code{force}
permits one.

\paragraph{Changed (B3b, REQ-47 and Section~\ref{sec:itceq}).} Three
author-seat calls at the B3b checkpoint, each raised by the role-review
passes. \stable{} REQ-47 now recognizes a reapplication from the
produced variables being present \emph{and} the History recording the
processor; matching names alone warn and then apply. Detection by names
alone refused a first application over a frame that legitimately
carried those names, which taught the caller to pass \code{force=True}
by reflex, the habit the requirement exists to prevent (DD-35). REQ-47
also admits \code{[meta] idempotent} as a declaration, so a
file-defined processor no longer has to be subclassed in Python to
declare that reapplying it is meaningful; it is the one typed field in
a strings-only section, and a quoted \code{"True"} is refused rather
than silently ignored (DD-36). Section~\ref{sec:itceq} gained three
section rules the parser enforces: a name declared in
\code{[constants]} may not also be an equation or correction target,
since a constant is substituted into every read and such an equation
would run with its result unreachable (DD-37); a forward reference in
file order is refused when the file is read, because the VarFrame may
carry the same name and the equation would silently use that value
before a later line overwrote it; and \code{[corrections]} may replace
a variable the VarFrame supplies, not only one \code{[equations]}
produced, which is the ordinary case of correcting a measured channel
in place. Requirement scope is otherwise unchanged.

\paragraph{Added (docs adoption).} Non-functional requirement REQ-108
(published documentation site on GitHub Pages, ProperDocs with markdown
SRS chapters and a build-time-generated API reference, mirroring the
pyflightstream site), draft, in Chapter~\ref{ch:nonfunctional}. A
pyflightstream adoption (DD-23); recorded in
\code{PYFLIGHTSTREAM\_ADOPTIONS.md}.

\section*{Document 0.1.1, 2026-07-21, M0 Uncertainty-Semantics Validation}

Issued at the M0 Phase 4 checkpoint, with the implementation and its
test suite as evidence.

\paragraph{Changed.} REQ-98 and REQ-99 promoted from draft to stable;
the smooth and diff row of REQ-98 stays provisional until OQ-18 is
resolved during v0.2.0 implementation, and operations without a frozen
weight rule raise when uncertainty is present. REQ-01 and REQ-03 gain
the dict-mode \code{"*"} swept sentinel (OQ-19). REQ-01 restricts
string-valued columns to dimension coordinates (OQ-21). REQ-11 states
the draft-guard scope: result exports guarded, inspection artifacts
exempt (OQ-22). Section~4.2 states that declared correlation applies
to both uncertainty components (OQ-23).

\paragraph{Added.} The \code{source\_coords} field of the
\code{Provenance} object, backing \code{db.manifest} (OQ-20).

\section*{Document 0.1.0, 2026-07-21, First Workspace-Tracked Version}

This version establishes the full requirement and design baseline of
\itc{} and starts requirement-evolution monitoring in the research
workspace. It consolidates the founding v0.1.0b draft with the 2026-07-21
specification review.

\paragraph{Added.} Functional requirements REQ-01 through REQ-73 and
REQ-98 through REQ-103; non-functional requirements REQ-74 through REQ-97;
non-requirements NREQ-01 through NREQ-10; design decisions DD-01 through
DD-22. The \code{.itc}, \code{.itceq}, \code{.itc\_pipe}, and
\code{.itc\_surr} file formats. The positioning statement among existing
tools (Section~\ref{sec:positioning}). The MIT license and citation policy
(Chapter~\ref{ch:contributing}). The incremental release roadmap mapping
milestones M0 through M3 to public releases \code{v0.1.0} through
\code{v0.4.0} (Chapter~\ref{ch:roadmap}).

\paragraph{Changed relative to the founding draft.} Two-component
uncertainty replaces the single-component UncFrame (REQ-99, DD-19);
uncertainty semantics are now normative for every operation (REQ-98,
DD-18); the axis system gains condition-dependent frames and moment
reference-point transfer (REQ-100, REQ-101); array-level immutability and
the state-hash scope are now requirements (REQ-102, REQ-103); expression
parsing is specified on the standard-library \code{ast} module (DD-20);
\code{ruff format} replaces \code{black}; the Python upper bound moves to
\code{<3.14}; the single-release plan is superseded by incremental
releases (DD-21).
